Rozrzut między najszybszą a~najwolniejszą bazą rośnie z~rozmiarem w~większości
przypadków, ale nie we wszystkich. Rośnie tam, gdzie co najmniej jedna baza
skanuje dane w~czasie $O(n)$, a~inna ma ścieżkę stałą lub logarytmiczną.
Przykłady to scenariusze bez indeksów R1, R2, R5, U3, D1, D2, D5 i~J2 oraz
R6, gdzie Valkey odczytuje licznik, a~pozostałe bazy agregują wszystkie bilety
(tab.~\ref{tab:h2}). Przy R4, D4 i~D6 bez indeksów rozrzut nie rośnie, bo
wszystkie bazy, także najszybsza, skanują tam pełne dane. Czas najszybszej
rośnie więc co najmniej w~tym samym tempie co czas najwolniejszej (Valkey ze
skanem \texttt{SCAN}).
Na przykład w~R1 bez indeksów MariaDB jest
\ratio{mariadb}{valkey}{500000}{off}{R1}$\times$ wolniejsza od Valkey przy
500 tys. biletów i~\ratio{mariadb}{valkey}{10000000}{off}{R1}$\times$ przy
10 mln.

Rozrzut nie rośnie dla operacji punktowych po kluczu (C1, C4, U1, U2, J3),
których czas w~każdej bazie jest stały. Nie rośnie też dla zapytań z~indeksem
o~małej selektywności, gdzie wszystkie bazy czytają kilka wierszy. H2 jest
\textbf{potwierdzona częściowo}. Różnice rosną z~liczbą rekordów wtedy, gdy
bazy mają różną złożoność ścieżki dostępu, czyli przy braku indeksu, przy
różnicy modeli danych (osadzenie a~złączenie) lub przy agregacji
utrzymywanej w~locie. Dla dobrze zaindeksowanych operacji punktowych rozmiar
zbioru nie ma znaczenia.
